% ==============================================================================
% Section 6: Conclusion
% ==============================================================================
\section{Conclusion}
\label{sec:conclusion}

This paper presents the first national-scale, outcome-disaggregated, monthly spatial analysis of disappearances in Mexico.
Covering \nmunis{} municipalities over \nmonths{} months (\studyperiod{}), it examines four outcome statuses---Total, Not Located, Located Alive, and Located Dead---through \nLISAruns{} LISA runs.

Five empirical facts are established.
First, disappearance outcome statuses are spatially non-interchangeable.
Jaccard similarity between Located Dead and either Not Located (0.143) or Located Alive (0.151) hotspots confirms that lethal-outcome geography is fundamentally distinct from the geography of active disappearances or living recoveries.
Second, geographic concentration exceeds standard criminological benchmarks by a factor of two to three: 50\% of Not Located cases in \referencemonth{} concentrated in just 42 municipalities (1.7\% of \nmunis{}), compared to the 4--6\% benchmark posited by \citet{weisburd_2015}.
Third, Centro is the fastest-growing hotspot region, with Not Located HH municipalities expanding at $+4.96$ per year ($p < 0.001$)---a trajectory invisible to annual analyses.
Fourth, the ``rising Baj\'{i}o'' hypothesis is rejected: Baj\'{i}o Total HH municipalities declined at $-1.09$ per year ($p = 0.024$).
Fifth, the CED's characterization of spatially and temporally heterogeneous attacks across Mexican states finds empirical support in the LISA dynamics, though with important caveats for states with low HH counts.

The policy implication is direct: outcome-specific geographic targeting is necessary.
Forensic identification resources should follow Located Dead hotspot geography, currently concentrated in Guanajuato and scattered municipalities in Tamaulipas and Jalisco.
Active search and prevention resources should follow Not Located hotspot geography, dominated by the rapidly expanding Centro corridor.
These are different places.
A composite ``disappearance map'' conflating all statuses would direct resources to a Located Alive--dominated geography that obscures the distinct municipalities where the most severe outcomes concentrate.

Future work will address the limitations of raw-count analysis.
A companion paper will introduce population-normalized rates and evaluate whether the non-interchangeability result holds under dual-metric (count and rate) analysis.
A subsequent econometric panel study will incorporate external covariates---organized crime incidence, prosecutorial capacity, socioeconomic indicators---to move from spatial description to causal inference.
